% Direct source weak-law proofs and an independent complete necessity argument.
\section{Weak laws and triangular arrays: proofs}
\label{sec:annex-p15-weak-laws-and-triangular-arrays}

The statements are in Chapter~\ref{sec:p15-weak-laws-and-triangular-arrays}.
The direct source extent is Durrett, Section 2.2, pp.~56--64
\cite{DurrettPTEAuthor2019}. None of these weak laws is proved by
invoking a strong law.

\begin{annexitem}{annex:p15.direct-variance-weak-laws}{thm:p15-direct-variance-weak-laws}{The full square expansion and the four conclusions}
\begin{proof}
Put \(\mu_i=\mathbb EX_i\). Finite second moments imply finite
first moments, and Cauchy--Schwarz makes every product of centered
variables integrable. Finite linearity gives
\(\mathbb ES_n=\sum_i\mu_i\), and expansion gives
\[
 \begin{aligned}
 \operatorname{Var}(S_n)
 &=\mathbb E\!\left[\left(\sum_i(X_i-\mu_i)\right)^2\right]\\
 &=\sum_i\mathbb E[(X_i-\mu_i)^2]
       +2\sum_{i<j}\mathbb E[(X_i-\mu_i)(X_j-\mu_j)].
 \end{aligned}
\]
For \(i\ne j\), the cross expectation is
\[
 \mathbb E(X_iX_j)-\mu_i\mathbb EX_j-\mu_j\mathbb EX_i+\mu_i\mu_j
       =\mathbb E(X_iX_j)-\mu_i\mu_j=0
\]
by the uncorrelatedness assumption. This proves the first assertion.
For any real \(c\), direct centering also gives
\(\operatorname{Var}(cY)=c^2\operatorname{Var}(Y)\).

For the second assertion, Markov's inequality applied to the
nonnegative \(|Z_n|^p\) at \(\varepsilon^p>0\) gives the displayed
bound with a non-strict threshold. The strict deviation event is
contained in it, so its probability also tends to zero.
For the common-mean sequence, the just-proved variance addition gives
\[
 \mathbb E[(S_n/n-\mu)^2]
   =n^{-2}\operatorname{Var}(S_n)
   =n^{-2}\sum_{i=1}^n\operatorname{Var}(X_i)\leq C/n.
\]
This proves mean-square convergence, and the second assertion with
\(p=2\) gives convergence in probability.
Finally, for arbitrary finite-second-moment \(S_n\),
\[
 \mathbb E\!\left[((S_n-\mu_n)/b_n)^2\right]=\sigma_n^2/b_n^2.
\]
The assumed limit is therefore mean-square convergence to zero,
and again the second assertion yields convergence in probability.
No relation between the probability spaces for different rows is
used in these deviation bounds.
\end{proof}
\end{annexitem}

\begin{annexitem}{annex:p15.bernstein-approximation}{thm:p15-bernstein-approximation}{The binomial representation and a uniform bound}
\begin{proof}
A continuous function on \([0,1]\) is bounded and uniformly
continuous. One elementary justification is sequential compactness:
successive bisection of \([0,1]\), always retaining a closed half
containing infinitely many terms, produces a convergent subsequence
of any sequence in the interval. An unbounded continuous function
would give a sequence with values of growing magnitude and then a
convergent subsequence contradicting continuity. Failure of uniform
continuity would give pairs \(u_j,v_j\) with \(|u_j-v_j|\to0\)
but function differences bounded away from zero; a convergent
subsequence of the \(u_j\) makes both arguments tend to the same
point, giving the same contradiction. Thus \(M_f<\infty\) and
\(\omega_f(\delta)\to0\) as \(\delta\downarrow0\).

Fix \(x\in[0,1]\), and take \(n\) independent Bernoulli variables,
each equal to one with probability \(x\) and zero with probability
\(1-x\). Their sum \(S\) has mass
\(\binom nm x^m(1-x)^{n-m}\) at \(m\): choose the \(m\) successful
positions and multiply the independent probabilities. Empty products
give the stated endpoint convention at \(x=0,1\).
Consequently \(f_n(x)=\mathbb E f(S/n)\),
\(\mathbb E(S/n)=x\), and
\(\operatorname{Var}(S/n)=x(1-x)/n\leq1/(4n)\).
For \(\delta>0\), split according to \(|S/n-x|\leq\delta\).
The integral absolute-value bound and Chebyshev give
\[
 \begin{aligned}
 |f_n(x)-f(x)|
 &\leq\mathbb E|f(S/n)-f(x)|\\
 &\leq\omega_f(\delta)+2M_f\mathbb P(|S/n-x|>\delta)\\
 &\leq\omega_f(\delta)+M_f/(2n\delta^2).
 \end{aligned}
\]
The bound is independent of \(x\), so it also bounds the supremum.
First choose \(\delta\) to make the modulus small, then choose \(n\)
large. This proves uniform convergence, rather than just convergence
at each fixed \(x\). The formula is a finite sum of polynomials of
degree at most \(n\).
For \(f(x)=x^2\), the expectation representation gives
\(f_n(x)=\operatorname{Var}(S/n)+[\mathbb E(S/n)]^2
=x^2+x(1-x)/n\). Its nonnegative error is maximized at \(x=1/2\),
where it equals \(1/(4n)\).
\end{proof}
\end{annexitem}

\begin{annexitem}{annex:p15.cube-shell-concentration}{thm:p15-cube-shell-concentration}{From independent coordinates to normalized volume}
\begin{proof}
Let \(U_1,\ldots,U_n\) be independent uniform coordinates on
\((-1,1)\). The law of their vector is normalized product Lebesgue
measure on the cube, whose volume is \(2^n\).
For \(Y_i=U_i^2\), direct integration gives
\(\mathbb EY_i=(1/2)\int_{-1}^1u^2\,du=1/3\) and
\(\mathbb EY_i^2=(1/2)\int_{-1}^1u^4\,du=1/5\).
Thus \(\operatorname{Var}(Y_i)=4/45\); the \(Y_i\) remain independent.
Putting \(M_n=n^{-1}\sum_iY_i\), the direct variance method gives
\(\mathbb EM_n=1/3\) and \(\operatorname{Var}(M_n)=4/(45n)\).

The radius relative to \(\sqrt{n/3}\) is \(R_n=\sqrt{3M_n}\).
For every nonnegative \(R\),
\(|R-1|=|R^2-1|/(R+1)\leq|R^2-1|\).
Outside the open annulus, \(|R_n-1|\geq\varepsilon\), so
\(|M_n-1/3|\geq\varepsilon/3\).
Chebyshev, including its non-strict threshold form, therefore bounds
the probability of the complement by
\[
 \frac{4}{45n}\frac{9}{\varepsilon^2}
       =\frac{4}{5n\varepsilon^2}.
\]
It is also at most one. Translating the probability through the
uniform product law proves the relative-volume assertion for every
\(\varepsilon>0\), including \(\varepsilon\geq1\).
\end{proof}
\end{annexitem}

\begin{annexitem}{annex:p15.coupon-collector-law}{thm:p15-coupon-collector-law}{Finite histories, geometric moments and coupon concentration}
\begin{proof}
Each event \(\{\tau_k^n\leq M\}\) depends on the first \(M\)
draws, so the extended hitting times and their finite-hitting event
are measurable. The specified finite replacements make the durations
and \(T_n\) real measurable variables.
For a fixed coupon label, the probability it is missing after \(M\)
draws is \((1-1/n)^M\), which tends to zero. Thus the probability
it never appears is zero. A union over the finitely many labels
proves \(\tau_n^n<\infty\) almost surely.

After a history ending at \(\tau_{k-1}^n=t\), exactly \(k-1\)
labels have been seen. The probability that the next \(d-1\) draws
belong to this set and the next draw is new equals
\[
 ((k-1)/n)^{d-1}(n-k+1)/n.
\]
This follows directly by independence of the future draws from the
fixed finite history, and is independent of the identities of the
seen labels or the history's length. To obtain joint duration
probabilities, prescribe \(D_{n,1}=d_1,\ldots,D_{n,k-1}=d_{k-1}\);
their last hitting index is the fixed number \(t=\sum_{j<k}d_j\).
Sum the preceding continuation probability times each prefix
probability over the disjoint finite histories realizing that duration
pattern. This multiplies the preceding joint probability by the same
geometric factor. Induction gives the product of all geometric masses.
Null exceptional paths assigned duration one do not change these
probabilities. This proves independence and the specified laws;
at \(k=1\) the first draw is always new, giving duration one.

For completeness, the geometric moment calculation used by the
source is supplied with its differentiation justified. Let \(0<p<1\)
and \(q=1-p\). A geometric \(G\) has masses \(pq^{d-1}\) for
\(d\geq1\), normalized by the geometric-series identity. Differentiate
the finite polynomial identity
\[
 \sum_{j=0}^Nq^j=(1-q)^{-1}-q^{N+1}(1-q)^{-1}.
\]
At fixed \(0<q<1\), the remainder and its first two derivatives tend
to zero: those derivatives are finite sums of terms with factors
bounded by a constant times \(N^2q^{N-1}\).
For each fixed integer \(d\geq0\), \(N^d q^N\to0\), since the
ratio of consecutive terms tends to \(q<1\) and is eventually
bounded by a number below one. Passing to the derivative limits gives
\[
 \sum_{j\geq1}j q^{j-1}=(1-q)^{-2},\qquad
 \sum_{j\geq2}j(j-1)q^{j-2}=2(1-q)^{-3}.
\]
Consequently
\(\mathbb EG=1/p\),
\(\mathbb E[G(G-1)]=2q/p^2\), and
\(\operatorname{Var}(G)=2q/p^2+1/p-1/p^2=q/p^2\).
At \(p=1\), \(G=1\) almost surely, giving the same moments.
This retains the geometric differentiation calculation of source
Example 1.6.14 \cite[p.~35]{DurrettPTEAuthor2019} without an
unproved exchange of an infinite derivative and a sum.

The independent durations sum to \(T_n\) almost surely. Reindexing
their success probabilities gives
\[
 \mathbb ET_n=n\sum_{j=1}^n j^{-1}=nH_n,\qquad
 \operatorname{Var}(T_n)
  =n^2\sum_{j=1}^n j^{-2}-nH_n\leq2n^2.
\]
The last bound follows from
\(\sum_{j\geq1}j^{-2}\leq1+\int_1^\infty x^{-2}\,dx=2\).
Decreasing-integrand comparisons also give
\(\log n\leq H_n\leq1+\log n\).
The variance criterion with \(b_n=n\log n\), \(n\geq2\),
now gives the centered weak limit, because the variance ratio
is at most \(2/(\log n)^2\). Since \(H_n/\log n\to1\),
adding the deterministic center gives the other limit.
The two numerical values at 365 are obtained by evaluating
\(365\log365\) and the finite sum \(365\sum_{j=1}^{365}1/j\);
they are approximately \(2153.462534\) and \(2364.646023\),
respectively. No finite-sample equality between these two quantities
is asserted.
\end{proof}
\end{annexitem}

\begin{annexitem}{annex:p15.permutation-cycle-law}{thm:p15-permutation-cycle-law}{The full sequential permutation construction and endpoint failure}
\begin{proof}
Generate the permutation by assigning one image at a time.
Start at domain label \(i_1=1\). At each step \(k=1,\ldots,n\), choose \(j_k\)
uniformly from the \(n-k+1\) range labels not yet used, and set
\(\pi(i_k)=j_k\). If \(j_k\) has not yet been used as a domain,
let \(i_{k+1}=j_k\). If the current cycle has closed and labels
remain, let \(i_{k+1}\) be the least unused domain label.
Stop after \(n\) assignments.

The finite choices can be made by independent uniform ranks among
the remaining ordered range labels. The algorithm never repeats a
domain or a range. For any fixed permutation, its images determine
one and only one possible sequence of such choices, recursively
through the rule for \(i_k\). That sequence has probability
\(\prod_{k=1}^n(n-k+1)^{-1}=1/n!\).
Thus the output is the uniform permutation.

Before choosing an image on a currently open cycle, its start label
has no assigned incoming image. Every other earlier domain in that
open chain already occurs in the used range, as do all labels in
closed cycles. Hence exactly one available range label, the current
start, closes the cycle. All other available choices continue to an
unused domain. At step \(k\), for every possible history,
the probability of closing is therefore \(1/(n-k+1)\).
Summing over histories with any prescribed preceding indicator
pattern multiplies that pattern's probability by this fixed closing
probability or its complement. Induction gives the product
probabilities of every indicator pattern, proving mutual independence.
This includes the final indicator, which is deterministically one.

Their Bernoulli means and variances give
\[
 \mathbb EK_n=\sum_{j=1}^n j^{-1}=H_n,\qquad
 v_n=\sum_{j=1}^n(j^{-1}-j^{-2})=H_n-\sum_{j=1}^n j^{-2}.
\]
For \(\varepsilon>0\),
\(v_n/(\log n)^{1+2\varepsilon}
\leq(1+\log n)/(\log n)^{1+2\varepsilon}\to0\).
The direct variance criterion proves the centered rate. Taking
\(\varepsilon=1/2\) and using \(H_n/\log n\to1\) proves the
cycle-count weak law.

To prove the source's asserted failure at exponent \(1/2\) without
a central limit theorem, center each independent indicator. Its
absolute value is at most one and their sum is \(K_n-H_n\).
Also \(v_n\geq\log n-2\). For sufficiently large \(n\),
\(v_n\geq1\) and \(v_n\geq(\log n)/2\). The bounded fourth-moment
lemma gives
\[
 \mathbb P(|K_n-H_n|>\sqrt{v_n/2})
       \geq\frac{v_n}{4(1+3v_n)}\geq1/16.
\]
Since \(\sqrt{v_n/2}\geq\sqrt{\log n}/2\), this gives the
stated fixed nonvanishing deviation probability.
Finally the displayed nine-label permutation sends
\(1\to3\to6\to1\), \(2\to9\to7\to5\to2\), and
\(4\to8\to4\). Concatenating those cycles closes them at
positions \(3,7,9\), verifying the source example and its index
convention.
\end{proof}
\end{annexitem}

\begin{annexitem}{annex:p15.occupancy-weak-law}{thm:p15-occupancy-weak-law}{The full covariance calculation and occupancy limit}
\begin{proof}
Let \(A_i\) mean that box \(i\) is empty. For \(n\geq2\),
independence of the labelled ball choices gives
\(\mathbb P(A_i)=q_n=(1-1/n)^r\), and for \(i\ne j\),
\(\mathbb P(A_i\cap A_j)=s_n=(1-2/n)^r\).
With no balls these are probabilities of the certain event,
so the empty-product convention applies, including \(n=2\).
Since \(N_n=\sum_i\mathbf1_{A_i}\), finite expansion yields
\[
 \mathbb EN_n=nq_n,\qquad
 \mathbb EN_n^2=nq_n+n(n-1)s_n,
\]
and subtracting \((nq_n)^2\) gives the asserted variance formula.
The diagonal and off-diagonal terms have both been retained.

For \(0\leq u<1\), integration of \(1/(1-t)\) between zero
and \(u\) gives
\(-u/(1-u)\leq\log(1-u)\leq-u\).
If \(r_n/n\to c<\infty\), these bounds show
\[
 r_n\log(1-1/n)\to-c,\qquad
 r_n\log(1-2/n)\to-2c,
\]
the second identity being used for \(n\geq3\).
Hence \(q_n\to e^{-c}\), \(s_n\to e^{-2c}\).
The exact source variance expression divided by \(n^2\) is
\[
 (1-1/n)(s_n-q_n^2)+(q_n-q_n^2)/n,
\]
which tends to zero. The direct variance criterion at \(b_n=n\)
then gives \(N_n/n-q_n\to0\) in probability, yielding the source
finite-\(c\) result.

There is also a uniform finite variance bound: for \(n\geq2\),
\(1-2/n\leq(1-1/n)^2\), and both bases are nonnegative.
Raising to the nonnegative integer \(r\) gives \(s_n\leq q_n^2\),
including \(r=0\). The off-diagonal covariance contribution is
therefore nonpositive, and
\(\operatorname{Var}(N_n)\leq nq_n(1-q_n)\leq n/4\).
If \(r_n/n\to\infty\), the logarithmic upper bound gives
\(q_n\leq e^{-r_n/n}\to0\); this same variance bound proves
the additional zero limit. With \(r_n=n\), the finite-\(c\)
case has \(c=1\). Dependence is visible already for two boxes
and one ball: each empty-box event has probability \(1/2\),
but their intersection is impossible.
\end{proof}
\end{annexitem}

\begin{annexitem}{annex:p15.triangular-array-weak-law}{thm:p15-triangular-array-weak-law}{The source two-error decomposition}
\begin{proof}
For each fixed \(n\), all truncated entries are bounded by \(b_n\)
and remain independent. Their sum has mean \(a_n\) and variance
equal to the sum of the truncated variances.
If \(S_n=\overline S_n\), a deviation of \(S_n\) from \(a_n\)
is the same deviation of \(\overline S_n\). Thus
\[
 \begin{aligned}
 \mathbb P(|S_n-a_n|>\varepsilon b_n)
 &\leq\mathbb P(S_n\ne\overline S_n)
       +\mathbb P(|\overline S_n-a_n|>\varepsilon b_n)\\
 &\leq\sum_{k=1}^n\mathbb P(|X_{n,k}|>b_n)
       +(\varepsilon b_n)^{-2}
                    \operatorname{Var}(\overline S_n)\\
 &\leq u_n+\varepsilon^{-2}b_n^{-2}
                         \sum_{k=1}^n\mathbb E\overline X_{n,k}^{\,2}
       =u_n+\varepsilon^{-2}w_n.
 \end{aligned}
\]
The second line uses the union bound and Chebyshev, and the third
uses variance addition and \(\operatorname{Var}(Y)\leq\mathbb EY^2\).
The assumed limits give convergence in probability.
Every step concerns one row and requires no independence from other
rows. The finite bound only requires \(b_n>0\); the source asymptotic
statement retains the additional condition \(b_n\to\infty\).
None of these inequalities asserts necessity of the two sufficient
truncation tests.
\end{proof}
\end{annexitem}

\begin{annexitem}{annex:p15.levy-symmetric-maximal}{thm:p15-levy-symmetric-maximal}{First crossing and the sign of the future sum}
\begin{proof}
For \(1\leq k\leq n\), let
\[
 A_k=\{|T_k|>t,\ |T_j|\leq t\text{ for every }j<k\}.
\]
These first-crossing events are disjoint and their union is the
partial-sum maximum event. Split \(A_k\) into
\(A_k^+=A_k\cap\{T_k>t\}\) and
\(A_k^-=A_k\cap\{T_k<-t\}\).
The future sum \(R_k=T_n-T_k\) is independent of both events.
It has a symmetric law: jointly negating the independent symmetric
future variables preserves their product law and negates their sum.
Consequently \(\mathbb P(R_k\geq0)=\mathbb P(R_k\leq0)\geq1/2\).
This includes \(R_n=0\).

On \(A_k^+\cap\{R_k\geq0\}\) the terminal sum exceeds \(t\),
and on \(A_k^-\cap\{R_k\leq0\}\) it is below \(-t\).
Independence therefore gives
\[
 \mathbb P(A_k\cap\{|T_n|>t\})
       \geq\tfrac12\mathbb P(A_k^+)+\tfrac12\mathbb P(A_k^-)
       =\tfrac12\mathbb P(A_k).
\]
Sum over the disjoint \(A_k\). This proves the first inequality.
For the non-strict version use the first crossing of \(|T_k|\geq t\),
split at \(T_k\geq t\) and \(T_k\leq-t\), and use the same
argument with terminal threshold \(\geq t\).
Finally \(Y_k=T_k-T_{k-1}\) implies
\(|Y_k|\leq2\max_{0\leq j\leq n}|T_j|\), where \(T_0=0\).
The increment event with threshold \(2t\) is therefore contained
in the partial-sum maximum event with threshold \(t\), proving
the second assertion. Only probabilities of events were used;
no expectation of the summands is required.
\end{proof}
\end{annexitem}

\begin{annexitem}{annex:p15.iid-weak-tail-criterion}{thm:p15-iid-weak-tail-criterion}{The complete direct sufficiency and independent necessity proofs}
\begin{proof}
Write \(W=|X_1|\). First assume \(x\mathbb P(W>x)\to0\).
Use the triangular array \(X_{n,k}=X_k\) with \(b_n=n\).
Its first truncation sum is \(n\mathbb P(W>n)\to0\).
Its second truncation sum is
\(n^{-1}\mathbb E[X_1^2\mathbf1_{\{W\leq n\}}]\).
The power-tail identity, applied to the bounded truncated variable,
gives
\[
 \begin{aligned}
 \mathbb E[X_1^2\mathbf1_{\{W\leq n\}}]
 &=\int_0^n2y\,[\mathbb P(W>y)-\mathbb P(W>n)]\,dy\\
 &\leq\int_0^n2y\,\mathbb P(W>y)\,dy .
 \end{aligned}
\]
The exact first equality retains the endpoint \(W=n\);
for \(y<n\), the relevant event is \(y<W\leq n\).
In particular the later bound is an inequality, not equality with
the untruncated tail integral.
Set \(g(y)=2y\mathbb P(W>y)\), with \(g(0)=0\).
This Borel function tends to zero at infinity and is globally bounded:
on a bounded interval \(g(y)\leq2y\), and beyond a sufficiently
large point it is at most one. The rescaled functions \(g(nx)\)
are therefore uniformly bounded on \([0,1]\), and tend to zero
there pointwise, including at zero. Affine Lebesgue substitution
and bounded convergence give
\[
 \frac1n\int_0^n g(y)\,dy=\int_0^1g(nx)\,dx\longrightarrow0.
\]
The substitution is the already proved
Theorem~\ref{thm:p15-affine-lebesgue}.
Thus both source array conditions hold. Its truncated row center
is \(a_n=n\mu_n\), proving \(S_n/n-\mu_n\to0\) in probability.

For necessity, suppose some deterministic real \(c_n\) satisfies
\(S_n/n-c_n\to0\) in probability on \((\Omega,\mathcal F,\mathbb P)\).
Construct the finite product probability space
\((\Omega^2,\mathcal F\otimes\mathcal F,\mathbb Q)\), where
\(\mathbb Q=\mathbb P\otimes\mathbb P\). Let
\(\widehat X_i(\omega,\omega')=X_i(\omega)\) and
\(\widetilde X_i(\omega,\omega')=X_i(\omega')\).
All members of the two copied sequences together are independent:
each finite family of constraints is a rectangle, whose probability
factors between its two coordinates, and within each coordinate the
original independence factors the remaining constraints.
Consequently \(Y_i=\widehat X_i-\widetilde X_i\) are independent
by the disjoint-group theorem and identically distributed.
Their laws are symmetric because interchanging the two product
coordinates preserves product probability and negates every difference.
No boundedness or moment assumption is imposed on these copies.
Write \(\widehat S_n=\sum_{i=1}^n\widehat X_i\) and
\(\widetilde S_n=\sum_{i=1}^n\widetilde X_i\).

Let \(T_n=\sum_{i=1}^nY_i=\widehat S_n-\widetilde S_n\).
For \(\eta>0\), the union bound gives
\[
 \mathbb Q(|T_n|/n>\eta)
 \leq2\mathbb P(|S_n/n-c_n|>\eta/2)\longrightarrow0.
\]
The common centers cancel. The symmetric maximal inequality now
implies, for every \(\eta>0\),
\[
 \mathbb Q\!\left(\max_{i\leq n}|Y_i|>\eta n\right)
       \leq2\mathbb Q(|T_n|>\eta n/2)\longrightarrow0.
\]
Put \(q_n=\mathbb Q(|Y_1|>\eta n)\). Independence makes the
left side \(1-(1-q_n)^n\), which is at least \(1-e^{-nq_n}\).
If \(nq_n\) were bounded away from zero along a subsequence,
this last quantity would be bounded away from zero there as well.
Thus \(n\mathbb Q(|Y_1|>\eta n)\to0\) for every \(\eta>0\).

Since the original variables are real, the copy satisfies
\(\mathbb Q(|\widetilde X_1|\leq n/2)\to1\), and this probability
is at least \(1/2\) for all sufficiently large \(n\).
Independence of the two copies and the triangle inequality give
\[
 \begin{aligned}
 \mathbb P(|X_1|>n)\,
      \mathbb Q(|\widetilde X_1|\leq n/2)
 &\leq\mathbb Q(|\widehat X_1-\widetilde X_1|>n/2).
 \end{aligned}
\]
Multiply by \(n\) and use the preceding limit at \(\eta=1/2\).
It follows that \(n\mathbb P(W>n)\to0\).
For real \(x\geq1\), put \(n=\lfloor x\rfloor\); then
\(x\mathbb P(W>x)\leq(n+1)\mathbb P(W>n)\to0\).
This completes an independent necessity proof of the full tail condition.

To characterize all centers, under that condition the already proved
sufficiency makes \(\mu_n\) work. If \(c_n-\mu_n\to0\), the
triangle inequality and a split at \(\varepsilon/2\) make \(c_n\)
work too. Conversely, if both centers work, then for every
\(\varepsilon>0\),
\[
 \mathbf1_{\{|c_n-\mu_n|>\varepsilon\}}
 \leq\mathbb P(|S_n/n-c_n|>\varepsilon/2)
       +\mathbb P(|S_n/n-\mu_n|>\varepsilon/2).
\]
The right side tends to zero. Since the left side is a deterministic
zero or one, it is eventually zero. This proves \(c_n-\mu_n\to0\).

Finally, for \(0<r<1\), choose \(M\geq1\) so that
\(\mathbb P(W>x)\leq1/x\) for \(x\geq M\), using the tail condition.
The power-tail integral then gives
\[
 \mathbb EW^r
 \leq M^r+r\int_M^\infty x^{r-2}\,dx
 =M^r+\frac{r}{1-r}M^{r-1}<\infty.
\]
The restriction \(0<r<1\) is required for this calculation.
\end{proof}
\end{annexitem}

\begin{annexitem}{annex:p15.direct-integrable-weak-law}{thm:p15-direct-integrable-weak-law}{The source's two dominated-convergence calculations}
\begin{proof}
Write \(W=|X_1|\) and \(\mu=\mathbb EX_1\).
For \(x>0\),
\[
 x\mathbb P(W>x)\leq\mathbb E[W\mathbf1_{\{W>x\}}]\longrightarrow0
\]
by dominated convergence with dominator \(W\).
The same theorem, applied to
\(X_1\mathbf1_{\{W\leq n\}}\), gives \(\mu_n\to\mu\).
The weak-tail criterion gives \(S_n/n-\mu_n\to0\) in probability.
For any \(\varepsilon>0\), eventually
\(|\mu_n-\mu|\leq\varepsilon/2\), so
\[
 \mathbb P(|S_n/n-\mu|>\varepsilon)
       \leq\mathbb P(|S_n/n-\mu_n|>\varepsilon/2)\longrightarrow0.
\]
This is the direct truncation proof of the integrable weak law.
\end{proof}
\end{annexitem}

\begin{annexitem}{annex:p15.weak-law-counterexamples}{thm:p15-weak-law-counterexamples}{Computing each truncation boundary and the infinite-mean weak example}
\begin{proof}
In the first array, the only nonzero entry is \(2n>n\), so it
is deleted. Thus \(a_n=w_n=0\), \(u_n=1\), and \(S_n/n=2\).
All entries are constants and are independent. In the second array,
the value of magnitude \(n\) is retained, so \(u_n=0\),
\(w_n=n^{-2}n^2=1\), \(a_n=0\), and \(S_n/n=\varepsilon\).
One random entry and the other constant-zero entries form an independent
row. In the third array all entries have magnitude one and mean zero.
They are retained, giving \(u_n=0\), \(w_n=n/n^2=1/n\) and
\(a_n=0\), but their sum is \(n\varepsilon\).
The identical nonconstant entries are dependent.
For the deterministic cancelling row, both nonzero entries are
deleted and \(u_n=2\), but the original and truncated sums are both
zero. These examples prove the asserted failures and show why the
sufficient tests are not a characterization.

For the symmetric example, the geometric identity and nonnegative
integration give
\[
 \sum_{j\geq1}\frac{q^j}{j}
   =\int_0^q\sum_{j\geq1}t^{j-1}\,dt
   =\int_0^q\frac{dt}{1-t}=-\log(1-q)\quad(0\leq q<1).
\]
At \(q=1/2\) this is \(\log2\), so the stated two-sided
masses sum to one. Each part expectation is
\((2\log2)^{-1}\sum_{j\geq1}1/j=\infty\).
For \(2^m\leq x<2^{m+1}\),
\[
 \mathbb P(|X_1|>x)
   =\frac1{\log2}\sum_{j\geq m+1}\frac1{j2^j}
   \leq\frac{2^{-m}}{(m+1)\log2}.
\]
Thus \(x\mathbb P(|X_1|>x)\leq2/[(m+1)\log2]\to0\).
The bounded truncation is symmetric at every cutoff \(n\), so
\(\mu_n=0\). The weak-tail criterion gives the claimed weak limit.
Its infinite absolute moment makes the threshold series at \(a_n=n\)
diverge by the integer-tail bound. Feller's already proved theorem
then gives \(\limsup_n|S_n|/n=\infty\) almost surely.
This demonstrates that concentration in probability can coexist with
unbounded almost-sure normalized excursions.

For the last positive dyadic law, at \(x=2^j\),
\(\mathbb P(X_1>x)=\sum_{k\geq j+1}2^{-k}=2^{-j}\).
The product is one and does not tend to zero.
Necessity in the weak-tail criterion therefore excludes every
deterministic centering sequence. The required independent sequences
with the specified discrete laws exist by the countable independent
sequence construction.
\end{proof}
\end{annexitem}
